% \todo[inline]{TODO: tirar esse foco no artigo anterior logo no começo}

Python has become one of the most widely adopted programming languages due to its flexibility and developer productivity gains~\cite{chun2001core,dyer-fse:2022,malloy-emse:2019,7042702,DBLP:conf/usenix/Rossum07}. However, to remain competitive and suitable for modern software development, Python has continually evolved to incorporate new features. This continuous evolution requires existing code-bases to progressively rejuvenate in order to benefit from newer language capabilities and avoid software aging effects ~\cite{Lucas:jsp2024,Lucas:emse2025,Overbey:regrowing,Pirkelbauer:sc-rejuvenation}. These modernization efforts often involve repository-wide transformations to adopt newer typing syntax, improve compatibility with recent Python versions, simplify type annotations, and leverage advances in static analysis tooling.

To identify representative modernization efforts for qualitative investigation, we build upon a previous large-scale repository mining study that analyzed the evolution history of 141 popular open-source Python projects hosted on GitHub~\cite{refminer}. That study detected more than 674,000 Type Hint-related source code transformations and identified commits containing extensive type-related modifications, providing a rich set of candidate cases for further investigation.

% The results of this previous study revealed a strong predominance of transformations associated with \ths. In addition, the mining process contributed to identifying a subset of large-scale modernization efforts, characterized by commits containing a large volume of automated or semi-automated type-related transformations. These commits are defined as the top-ranked commits in terms of detected refactoring operations, which represent the most intensive modernization efforts observed in the dataset.

% Although the previous study provides a comprehensive quantitative characterization of Python modernization through large-scale mining, it does not explain the underlying decision-making processes that lead developers to perform such extensive transformations. In particular, it does not capture the motivations behind these modifications, the contextual factors influencing modernization decisions, nor the role of developers in orchestrating or validating automated refactoring activities.

To better understand the motivations for those large-scale modernization efforts, we built a curated dataset of 117 representative large-scale commits from an initial set containing the 200 most important commits from each developer. The curation process prioritized commits that best exemplified substantial source code transformations and provided sufficient contextual information for qualitative investigation. A detailed description of this selection procedure is presented in \ref{sec:approach}. Each selected commit was authored by a distinct developer which was used for our subsequent qualitative analysis and developer survey.

Understanding these aspects is essential to complement mining-based studies with evidence that captures developers rationale and the context surrounding modernization efforts. Therefore, in this work, we extend the previous study by investigating the motivations and practices related to large-scale Python modernization.

For this end, we conducted a mixed-method empirical study combining a developer survey with qualitative analysis and triangulation. The survey focused on (i) motivations for performing large-scale \ths-related transformations, (ii) tools and automation techniques used during the process, and (iii) perceptions regarding the potential use of Foundation Models to support software modernization tasks. We analyzed the survey responses using thematic analysis and triangulated the resulting themes with repository artifacts, including commit messages, pull request discussions, and code review comments. This triangulation enabled us to assess the consistency between developer-reported motivations and the evidence observed in the repository history, providing a more comprehensive understanding of large-scale modernization practices.

Overall, this work bridges the gap between large-scale repository mining and a developer-centered understanding of software modernization. Some of the main contributions are:

\begin{itemize}
    \item \textbf{A mixed-method investigation of large-scale Type Hint modernization.} We combine repository mining, developer surveys, thematic analysis, and qualitative triangulation to study the motivations, practices, and perceptions associated with large-scale modernization efforts.

    \item \textbf{An understanding of practices on Type Hints modernization.} We identify the technical and ecosystem-related motivations behind modernization campaigns and characterize the tools and practices developers adopt to perform them.

    \item \textbf{Insights into the emerging role of Foundation Models in software modernization.} We characterize developers' current perceptions of AI-assisted modernization, highlighting both growing acceptance and continued preference for deterministic tooling in rule-based transformations.

    \item \textbf{Evidence supporting the use of repository artifacts as a source of developer rationale.} Through qualitative triangulation, we show that repository artifacts and developer reports are largely consistent, while surveys contribute additional contextual explanations that are often absent from project history.
\end{itemize}

The remainder of this paper is organized as follows. Section \ref{sec:background} provides the background on Python language features and source-code transformations. Section \ref{sec:related-work} reviews the related work. Section \ref{sec:approach} describes the research approach, including the study design, data collection, and analysis procedures. Section \ref{sec:results} presents the results of the empirical study. Section \ref{sec:discussion} discusses the main findings and their implications. Section \ref{sec:threats} outlines the threats to the validity of the study. Finally, Section \ref{sec:conclusion} concludes the paper and discusses directions for future work.

% Python has become one of the most widely adopted programming languages due to its flexibility and developer productivity gains~\cite{chun2001core,dyer-fse:2022,malloy-emse:2019,7042702,DBLP:conf/usenix/Rossum07}. However, to remain competitive and suitable for modern software development, Python, like other programming languages, has continually evolved to incorporate modern features. This continuous evolution requires existing codebases to progressively rejuvenate in order to benefit from newer language capabilities and avoid software aging effects ~\cite{Lucas:jsp2024,Lucas:emse2025,Overbey:regrowing,Pirkelbauer:sc-rejuvenation}.

% Among the most relevant modernization efforts introduced into Python are \ths ~\cite{pep484} and \spm ~\cite{pep634}. The Type Hints, enable developers to incrementally annotate variables, parameters, and return types, improving static analysis, documentation quality, IDE support, and defect detection ~\cite{Khan:type-related-defects,10.1145/3540250.3549114}.  Structural Pattern Matching, extends the language with declarative pattern-based control flow using the \texttt{match/case} construct, allowing developers to replace traditional \texttt{if-elif-else} chains with more expressive and maintainable structures.  Together, these features represent important steps toward modernizing Python systems and improving software comprehensibility and maintainability over time.

% Although previous studies have investigated the adoption and benefits of Type Hints and Structural Pattern Matching ~\cite{10.1145/3540250.3549114,Khan:type-related-defects,10006855,10589769}, important aspects of their evolutionary adoption process remain underexplored. Existing research has mainly focused on static snapshots of adoption or code quality benefits, with limited attention given to how developers modernize legacy Python systems over time.

% To address these gaps, this paper investigate software rejuvenation practices in Python projects through the analysis of source-code transformations related to \ths and \spm. To achieve this goal, we extended the RefactoringMiner infrastructure ~\cite{refminer} and mined the evolution history of 141 popular open-source Python projects hosted on GitHub, collecting more than 674,000 transformation instances. Some of our
% findings are:
